\begin{corollary}[Predecessor Existence and Uniqueness Away from $1$]
\label{cor:predecessor-exists-unique-away-from-one}
Let $(P,S,1)$ be a Peano system, and let $x \in P$ with $x \neq 1$. Then
there exists a unique $u \in P$ such that $S(u)=x$.
\hyperref[prf:predecessor-exists-unique-away-from-one]{Go to proof.}
\end{corollary}
\end{corollarybox}
\LeanFormalizes{cor:predecessor-exists-unique-away-from-one}{lra-lean}{LRA.VolumeII.PeanoSystems.PeanoSystem.Predecessor}{PredecessorExistsUniqueAwayFromOne}{pending}

\begin{remark*}[Standard quantified statement]
For a Peano system $(P,S,1)$,
\[
\forall x \in P\;\bigl(x \neq 1 \Longrightarrow \exists! u \in P\;(S(u)=x)\bigr).
\]
\end{remark*}

\begin{remark*}[Predicate reading]
\[
\forall x \in P\;\bigl(x \neq 1 \Longrightarrow \exists! u \in P\;(S(u)=x)\bigr).
\]
\end{remark*}

\begin{remark*}[Interpretation]
Every element distinct from $1$ has exactly one predecessor. The successor map
therefore organizes the non-initial part of a Peano system into one-step
extensions of earlier elements.
\end{remark*}

\begin{remark*}[Source comparison]
Feferman proves predecessor existence and uniqueness together inside Theorem
3.3. The present notes separate the uniqueness statement into an independent
corollary so that it can be cited directly in later proofs.
\citet{FefermanNumberSystems1964}
\end{remark*}

\begin{dependencies}
\begin{itemize}
  \item \hyperref[def:peano-system]{Peano System}
  \item \hyperref[cor:non-one-elements-have-a-predecessor]{Non-One Elements Have a Predecessor}
  \item \hyperref[ax:peano-successor-injective]{Injectivity of Successor}
\end{itemize}
\end{dependencies}

